% Experimental results (V6 revision).
% Directive unchanged from V5: FULL graphical + tabular treatment for Phase B1
% and Phase C2 ONLY; every other phase is a short focused paragraph with no
% tables and no figures. V6 deepens the two focus phases (drop-vs-clean table,
% per-level margin table, L1-L5 sample strip, full Spearman table), weaves the
% plain-language reading into the prose, and corrects the V5 claims that
% contradicted the frozen data (ECE at L5, TransNeXt latency, the "under
% 1.5 pp from L3" margin). Phase C2 axis comparisons are plotted against
% measured PSNR and SSIM -- never the preset L1-L5 index.
\section{Experimental Results}
\label{sec:tests}

This section reports the measured performance. Every value comes from a
frozen training run recorded under \texttt{runs/final/} and aggregated by
\texttt{artifacts/Final\_Exp.json} (schema version 3); no values are
synthesized or interpolated. The two phases that carry the headline science
--- Phase~B1 (combined degradation) and Phase~C2 (single-axis isolation under
the THz protocol) --- are reported in full, with tables, figures, and a
model-by-model reading. The remaining phases are summarized in
Section~\ref{sec:tests:brief}; their complete per-cell numbers live in the
machine-readable artifacts referenced in the Appendix.

\subsection{Test rig}
\label{sec:tests:rig}
The campaign ran on a single workstation with the configuration in
Table~\ref{tab:rig}. Determinism is guarded by
\texttt{pl.seed\_everything(42, workers=True)} for the canonical seed and by
the per-sample seed offset for the validation split, so any cell can be
reproduced bit-exactly.

\begin{table}[htbp]
\centering
\caption{Test rig and software stack.}
\label{tab:rig}
\setlength{\tabcolsep}{4pt}
\renewcommand{\arraystretch}{1.1}
\begin{tabular}{ll}
\toprule
Component & Value \\
\midrule
GPU                & NVIDIA RTX 5070 (Blackwell, sm\_120, 12 GB) \\
Precision          & \texttt{bf16-mixed} \\
CPU dataloader workers & 0 (Windows) \\
Framework          & PyTorch 2.11 [17] + Lightning 2.x [16] \\
Image quality      & \texttt{torchmetrics} PSNR / SSIM, \texttt{data\_range = 1} \\
Determinism gate   & \texttt{src/tests/test\_degradation\_determinism.py} \\
\bottomrule
\end{tabular}
\end{table}

\subsection{What the degradations look like}
\label{sec:tests:samples}
Figure~\ref{fig:samples} grounds the numbers visually. Each grid shows the
same six fixed validation images under the clean pipeline, the combined
Phase~B1 recipe at the moderate level, the grayscale Phase~B2 protocol, and
each of the three Phase~C2 single-axis pipelines. Two things are worth
noticing before any accuracy number. First, the combined B1~L3 image is
already hard for a human: the object outline survives, but texture and color
are essentially gone. Second, the three C2 rows differ far less in
\emph{appearance} than in what they destroy --- the salt-and-pepper row looks
the most damaged but keeps the underlying shape intact, while the resolution
row looks smooth but has permanently discarded the fine structure. The
accuracy results below follow this asymmetry exactly.

\begin{figure}[htbp]
  \centering
  \includegraphics[width=0.49\textwidth]{\figpath{samples/grid_THz_L3_cifar10.png}}\hfill
  \includegraphics[width=0.49\textwidth]{\figpath{samples/grid_THz_L3_mnist.png}}
  \caption{Degradation pipelines at the moderate level (L3) on six fixed validation images. Left: CIFAR-10; right: MNIST. Rows: clean; combined Phase B1; grayscale Phase B2; the three Phase C2 single axes.}
  \label{fig:samples}
\end{figure}

\subsection{The remaining phases in brief}
\label{sec:tests:brief}
The five phases outside the B1 / C2 focus are summarized here; each paragraph
states what was run, the one or two numbers that matter, and what they mean.

\sectheading{Clean baselines (Phase A, 6 cells).} All three backbones
classify the undegraded validation sets at near-ceiling accuracy: 95.18\%
(ResNet50), 93.56\% (DenseNet121), and 97.64\% (TransNeXt-tiny) on CIFAR-10,
and 99.1--99.3\% on MNIST. These six numbers are the reference points every
drop in this section is measured against.

\sectheading{Full single-axis isolation (Phase C, 150 cells).} Sweeping all
five degradation axes one at a time, over the full color pipeline, produces a
clear severity ordering: resolution loss dominates, blur and additive noise
are intermediate, and desaturation and salt-and-pepper are comparatively
benign in isolation. In plain terms: taking away the pixels hurts far more
than taking away the color. This ordering is what motivated the three-axis
THz-relevant subset studied in depth as Phase~C2.

\sectheading{Regularization recovery (Phase D, 90 cells).} Layering dropout,
mixup, and CutMix on top of the combined-degradation pipeline recovers very
little: the best treatment (T3, the combination) recovers +0.77~\Mpp{} on
average against a collapse roughly sixteen times larger, with the best single
segment (TransNeXt-tiny on CIFAR-10) at +2.20~\Mpp{}. The conclusion is
important precisely because the effect is small: if the accuracy loss were
over-fitting, regularization would have clawed a large part of it back. It
does not, so the loss is an information bottleneck --- the detail the models
need is no longer in the pixels.

\sectheading{THz protocol simplification (Phase B2 / B2-nr, 36 cells).}
Re-running the combined sweep under the deployment-realistic protocol
(grayscale fixed, additive Gaussian noise off, T3 active) lifts accuracy by
+4.01~\Mpp{} on average over Phase~B1, and by +5.58~\Mpp{} at the moderate
level. The six-cell no-regularization arm (B2-nr) decomposes the L3 gain:
+4.19~\Mpp{} comes from the simpler protocol itself and only +1.39~\Mpp{}
from T3. In other words, about three quarters of the improvement is the
physics of the protocol, not the training trick --- consistent with the
Phase~D reading. A secondary observation: T3 helps roughly twice as much
under the simpler protocol (+1.39) as under the full one (+0.77), so
regularization becomes more useful as the degradation becomes less severe.

\sectheading{Multi-seed audit (48 replicates).} Re-training all 24 L3
headline cells at seeds 43 and 44 leaves every conclusion intact: the largest
per-cell standard deviation is 1.18~\Mpp{} and no cross-model ordering
reverses under any seed. Every margin quoted in this section is therefore
either comfortably above the seed noise or explicitly flagged as within it.

\sectheading{Calibration.} Expected calibration error stays moderate across
the campaign: the mean over the 24 audited cells is 5.6\% at L3 and 2.9\% at
L5 (L3 measured on the held-out test split, L5 on the validation split).
Two useful readings: the networks do not become wildly over-confident as the
input degrades --- confidence broadly tracks the collapsing accuracy --- and
the single worst miscalibration (12.7\%) occurs on the unregularized B2-nr
arm, hinting that the Phase~D regularization also buys calibration, not only
accuracy.

\sectheading{Inference throughput.} On the RTX 5070 at batch size 1, mean
per-image latency is 4.2~ms for ResNet50 (238 images/s), 10.7~ms for
DenseNet121 (94 images/s), and 27.5~ms for TransNeXt-tiny (36 images/s).
The transformer is thus roughly $6.5\times$ slower than ResNet50 --- its
foveal-attention path is not graph-compiled on the current stack --- but at
36 frames per second it still clears the acquisition rate of practical THz
scanners, which is the budget that matters for deployment.

\sectheading{Confusion structure.} At the extreme level the residual errors
concentrate among visually similar classes: cat--dog and bird--deer on
CIFAR-10, 4--7 and 8--3 on MNIST. Even after the fine detail is destroyed,
the networks keep a coarse notion of category; what they lose is the ability
to separate lookalikes.

\subsection{Phase B1: combined degradation}
\label{sec:tests:phase_b1}
Phase B1 applies all five degradation axes simultaneously at a common level
and asks the deployment-shaped question: how much accuracy survives when
everything degrades at once? Here the L1--L5 index is a legitimate $x$-axis,
because Phase~B1 sweeps a single fixed recipe and each level is one
well-defined operating point. Figure~\ref{fig:b1_samples} shows what the
sweep does to the images; Table~\ref{tab:phase_b} reports the full
$3 \times 2 \times 5$ accuracy block; Figure~\ref{fig:b1_curves} plots the
curves against the per-model clean baselines; and
Tables~\ref{tab:b1_delta}--\ref{tab:b1_gap} quantify the two readings the
curves suggest --- how far each model falls, and who leads by how much.

\begin{figure}[htbp]
  \centering
  \includegraphics[width=0.49\textwidth]{\figpath{samples/grid_B1_levels_cifar10.png}}\hfill
  \includegraphics[width=0.49\textwidth]{\figpath{samples/grid_B1_levels_mnist.png}}
  \caption{The combined Phase B1 recipe at L1--L5 on six fixed validation images. Left: CIFAR-10; right: MNIST.}
  \label{fig:b1_samples}
\end{figure}

\begin{table}[htbp]
\centering
\caption{Phase B (combined degradation) validation accuracy (\%). All five axes active at the same level L. Cells with $\pm$ values are multi-seed audited (seeds $\{42, 43, 44\}$).}
\label{tab:phase_b}
\setlength{\tabcolsep}{4pt}
\renewcommand{\arraystretch}{1.1}
\begin{tabular}{llrrrrr}
\toprule
Model & Dataset & L1 & L2 & L3 & L4 & L5 \\
\midrule
ResNet50 & CIFAR-10 & 76.46 & 53.94 & 36.31 $\pm$ 0.83 & 26.84 & 19.06 \\
ResNet50 & MNIST & 99.18 & 96.36 & 79.17 $\pm$ 0.80 & 56.44 & 28.00 \\
DenseNet121 & CIFAR-10 & 78.48 & 53.68 & 38.54 $\pm$ 0.78 & 28.52 & 19.58 \\
DenseNet121 & MNIST & 98.74 & 96.54 & 80.79 $\pm$ 0.84 & 58.08 & 27.76 \\
TransNeXt-tiny & CIFAR-10 & 87.58 & 64.22 & 39.48 $\pm$ 0.62 & 28.50 & 20.10 \\
TransNeXt-tiny & MNIST & 98.96 & 96.00 & 80.05 $\pm$ 0.11 & 57.24 & 27.48 \\
\bottomrule
\end{tabular}
\end{table}


\begin{figure}[htbp]
  \centering
  \includegraphics[width=\textwidth]{\figpath{curves/b1_acc_vs_level.png}}
  \caption{Phase B1 validation accuracy versus combined level; dotted lines mark the per-model clean baselines. Left: CIFAR-10; right: MNIST.}
  \label{fig:b1_curves}
\end{figure}

\begin{table}[htbp]
\centering
\caption{Phase B1 accuracy drop versus the Phase A clean baseline (percentage points); the Clean column is the Phase A accuracy (\%). L3 uses the multi-seed mean.}
\label{tab:b1_delta}
\setlength{\tabcolsep}{4pt}
\renewcommand{\arraystretch}{1.1}
\begin{tabular}{llrrrrrr}
\toprule
Model & Dataset & Clean & $\Delta$L1 & $\Delta$L2 & $\Delta$L3 & $\Delta$L4 & $\Delta$L5 \\
\midrule
ResNet50 & CIFAR-10 & 95.18 & -18.72 & -41.24 & -58.87 & -68.34 & -76.12 \\
ResNet50 & MNIST & 99.14 & +0.04 & -2.78 & -19.97 & -42.70 & -71.14 \\
DenseNet121 & CIFAR-10 & 93.56 & -15.08 & -39.88 & -55.02 & -65.04 & -73.98 \\
DenseNet121 & MNIST & 99.30 & -0.56 & -2.76 & -18.51 & -41.22 & -71.54 \\
TransNeXt-tiny & CIFAR-10 & 97.64 & -10.06 & -33.42 & -58.16 & -69.14 & -77.54 \\
TransNeXt-tiny & MNIST & 99.24 & -0.28 & -3.24 & -19.19 & -42.00 & -71.76 \\
\bottomrule
\end{tabular}
\end{table}


\begin{table}[htbp]
\centering
\caption{TransNeXt-tiny accuracy margin over the two CNNs on the Phase B1 curve (percentage points); L3 uses the multi-seed mean (seed noise $\sigma < 1.2$).}
\label{tab:b1_gap}
\setlength{\tabcolsep}{5pt}
\renewcommand{\arraystretch}{1.1}
\begin{tabular}{llrrrrr}
\toprule
Dataset & Margin & L1 & L2 & L3 & L4 & L5 \\
\midrule
CIFAR-10 & vs CNN mean & +10.11 & +10.41 & +2.05 & +0.82 & +0.78 \\
CIFAR-10 & vs best CNN & +9.10 & +10.28 & +0.94 & -0.02 & +0.52 \\
MNIST & vs CNN mean & +0.00 & -0.45 & +0.07 & -0.02 & -0.40 \\
MNIST & vs best CNN & -0.22 & -0.54 & -0.75 & -0.84 & -0.52 \\
\bottomrule
\end{tabular}
\end{table}


Three observations organize the block.

\textbf{The collapse is monotonic and severe on CIFAR-10.}\quad Accuracy
falls with every level step for every model
(Table~\ref{tab:b1_delta}): even the mild level already costs 10--19~\Mpp{},
the moderate level costs 54--59, and by L5 all three models have lost about
three quarters of their clean accuracy and sit at 19--21\% --- better than
the 10\% random-guess floor, but not by much. Nothing in the design space
escapes this fall; the question the later phases answer is \emph{why} (it is
the information, not the training --- Phases~D and C2).

\textbf{The transformer leads exactly where the images are still
readable.}\quad Table~\ref{tab:b1_gap} makes the margin structure explicit.
TransNeXt-tiny leads the CNN mean by +10.11~\Mpp{} at L1 and +10.41 at L2 ---
an advantage far above the seed noise --- but the lead collapses to
+2.05~\Mpp{} at L3 (only +0.94 over the better CNN, within about one
combined seed-$\sigma$), and to under 1~\Mpp{} at L4--L5. Comparing with
Figure~\ref{fig:b1_samples}: the transformer's foveal attention pays off
while local texture still exists (L1--L2), and stops paying the moment the
degradation has destroyed the texture it attends to. From L3 onward the
three architectures are practically interchangeable.

\textbf{MNIST resists four levels, then falls off the same cliff.}\quad The
digit task barely notices the first two levels (drops under 3.3~\Mpp{};
TransNeXt has no advantage there --- the CNNs are already at ceiling), stays
usable at L3 (about 79--81\%), and then collapses to 27--28\% at L5, a loss
of about 71~\Mpp{} for every model. High-contrast shapes on a black
background simply survive the pipeline longer than natural images; but once
downsampling to $3 \times 3$ pixels removes the stroke structure, MNIST
converges to the same information floor as CIFAR-10. The per-model reading
of both datasets continues in Section~\ref{sec:crossmodel}.

\subsection{Phase C2: single-axis attribution under the THz protocol}
\label{sec:tests:phase_c2}
Phase C2 answers the engineering question behind the whole campaign: of the
degradations a real THz sensor produces, which one should be fought first?
Each cell isolates one of the three THz-relevant axes --- resolution, blur,
or salt-and-pepper --- at one level, with the other two held at identity,
under the THz protocol (image fixed at full grayscale, additive Gaussian
noise off, T3 regularization). Table~\ref{tab:phase_c2} reports the full
accuracy block.

\begin{table}[htbp]
\centering
\caption{Phase C2 single-axis isolation under the THz protocol (saturation $=0$, noise $=0$, T3 regularization) validation accuracy (\%).}
\label{tab:phase_c2}
\setlength{\tabcolsep}{4pt}
\renewcommand{\arraystretch}{1.1}
\begin{tabular}{lllrrrrr}
\toprule
Model & Dataset & Axis & L1 & L2 & L3 & L4 & L5 \\
\midrule
ResNet50 & CIFAR-10 & Resolution & 80.18 & 66.36 & 51.32 & 42.96 & 30.16 \\
ResNet50 & CIFAR-10 & Blur & 90.62 & 89.58 & 83.36 & 71.68 & 58.76 \\
ResNet50 & CIFAR-10 & S\&P & 86.24 & 85.18 & 83.48 & 77.86 & 74.68 \\
ResNet50 & MNIST & Resolution & 99.06 & 98.64 & 94.12 & 81.24 & 45.54 \\
ResNet50 & MNIST & Blur & 99.24 & 99.26 & 99.32 & 99.32 & 98.86 \\
ResNet50 & MNIST & S\&P & 99.14 & 99.16 & 99.18 & 99.00 & 99.00 \\
DenseNet121 & CIFAR-10 & Resolution & 76.82 & 64.04 & 51.34 & 44.44 & 30.80 \\
DenseNet121 & CIFAR-10 & Blur & 88.24 & 85.66 & 81.66 & 74.90 & 64.60 \\
DenseNet121 & CIFAR-10 & S\&P & 86.38 & 83.26 & 80.54 & 79.30 & 77.78 \\
DenseNet121 & MNIST & Resolution & 99.38 & 98.64 & 94.30 & 82.28 & 45.02 \\
DenseNet121 & MNIST & Blur & 99.42 & 99.34 & 99.22 & 99.38 & 99.16 \\
DenseNet121 & MNIST & S\&P & 99.16 & 99.20 & 99.20 & 99.08 & 99.06 \\
TransNeXt-tiny & CIFAR-10 & Resolution & 89.34 & 76.08 & 58.24 & 46.58 & 30.54 \\
TransNeXt-tiny & CIFAR-10 & Blur & 96.52 & 95.96 & 93.00 & 88.02 & 71.04 \\
TransNeXt-tiny & CIFAR-10 & S\&P & 95.20 & 94.78 & 93.52 & 91.56 & 89.58 \\
TransNeXt-tiny & MNIST & Resolution & 99.32 & 98.76 & 93.94 & 81.74 & 45.60 \\
TransNeXt-tiny & MNIST & Blur & 99.38 & 99.38 & 99.46 & 99.48 & 99.06 \\
TransNeXt-tiny & MNIST & S\&P & 99.28 & 99.14 & 99.20 & 99.08 & 99.08 \\
\bottomrule
\end{tabular}
\end{table}


\textbf{Why the plots use PSNR and SSIM, not the level index.}\quad A given
nominal level does not do equal damage on different axes: L3 blur removes far
less information than L3 resolution, so comparing axes at matched
\emph{level} would compare unmatched images. The level is a knob we set; PSNR
and SSIM are what the images actually measure. Every Phase~C2 comparison is
therefore drawn against the measured image quality --- PSNR
(Figure~\ref{fig:c2_type_psnr}) and SSIM (Figure~\ref{fig:c2_type_ssim})
separately --- computed against the clean upsampled reference, with the
L1/L5 endpoints of the resolution curve annotated for orientation.

\begin{figure}[htbp]
  \centering
  \includegraphics[width=\textwidth]{\figpath{curves/c2_type_psnr.png}}
  \caption{Phase C2 axis comparison versus measured PSNR, one curve per axis averaged over the backbones; L1/L5 mark the resolution-axis endpoints. Left: CIFAR-10; right: MNIST.}
  \label{fig:c2_type_psnr}
\end{figure}

\begin{figure}[htbp]
  \centering
  \includegraphics[width=\textwidth]{\figpath{curves/c2_type_ssim.png}}
  \caption{Phase C2 axis comparison versus measured SSIM; same construction as Figure~\ref{fig:c2_type_psnr}.}
  \label{fig:c2_type_ssim}
\end{figure}

\textbf{Resolution is the axis that matters.}\quad At every matched PSNR and
SSIM value, the resolution curve sits far below blur, and blur below
salt-and-pepper; the ordering is identical under both quality metrics. Averaged
over the six (model, dataset) pairs, the L1--L5 mean accuracy cost is
$-27.45$~\Mpp{} for resolution against $-6.58$ for blur and $-5.13$ for
salt-and-pepper --- a factor of roughly four. The plain-language version of
Figure~\ref{fig:c2_type_psnr} is: an image that \emph{measures} equally
degraded is much more dangerous when the damage came from lost pixels than
when it came from smoothing or speckle, because only the former removes
information the classifier cannot re-learn its way around.

\begin{table}[htbp]
\centering
\caption{Phase C2: Spearman rank correlation $\rho$ between accuracy and measured image quality, pooled over the three backbones ($n = 15$ per entry).}
\label{tab:c2_spearman}
\setlength{\tabcolsep}{6pt}
\renewcommand{\arraystretch}{1.1}
\begin{tabular}{lrrrr}
\toprule
 & \multicolumn{2}{c}{CIFAR-10} & \multicolumn{2}{c}{MNIST} \\
Axis & $\rho$ vs PSNR & $\rho$ vs SSIM & $\rho$ vs PSNR & $\rho$ vs SSIM \\
\midrule
Resolution & +0.98 & +0.98 & +0.98 & +0.98 \\
Blur & +0.84 & +0.84 & +0.42 & +0.42 \\
Salt \& pepper & +0.53 & +0.53 & +0.70 & +0.70 \\
\bottomrule
\end{tabular}
\end{table}


\textbf{The quality metrics rank the axes but do not explain them.}\quad
Table~\ref{tab:c2_spearman} quantifies how tightly accuracy follows measured
quality. Pooled across models, resolution shows an almost perfect
association ($\rho = +0.98$ on both datasets and both metrics); blur and
salt-and-pepper are markedly weaker ($+0.84$ and $+0.53$ on CIFAR-10). Two
nuances make this table honest rather than decorative. First, within any
single model on CIFAR-10 the association is perfectly monotone ($\rho = 1.0$)
on every axis --- the pooled values drop below 1 only because the three
models sit at different accuracy offsets at the same quality. Second, on
MNIST the blur and salt-and-pepper correlations are weak for the opposite
reason: those axes barely move MNIST accuracy at all (the curves are flat
near 99\%), so there is almost no signal for a rank correlation to find. A
low $\rho$ here means ``harmless axis,'' not ``unpredictable axis.''

\begin{figure}[htbp]
  \centering
  \includegraphics[width=\textwidth]{\figpath{curves/c2_model_psnr.png}}
  \caption{Phase C2 model comparison versus PSNR: one panel per axis (columns) and dataset (rows), one curve per backbone.}
  \label{fig:c2_model_psnr}
\end{figure}

\begin{figure}[htbp]
  \centering
  \includegraphics[width=\textwidth]{\figpath{curves/c2_model_ssim.png}}
  \caption{Phase C2 model comparison versus SSIM; same construction as Figure~\ref{fig:c2_model_psnr}.}
  \label{fig:c2_model_ssim}
\end{figure}

\textbf{Model ordering under the THz protocol.}\quad
Figures~\ref{fig:c2_model_psnr} and \ref{fig:c2_model_ssim} repeat the
comparison with one curve per model. On CIFAR-10, TransNeXt-tiny is the most
robust backbone on every axis across the entire quality range --- most
visibly on blur, where at L5 it holds 71.0\% against 58.8\% (ResNet50) and
64.6\% (DenseNet121) --- and the two CNNs essentially overlap. On MNIST only
the resolution axis produces a meaningful fall (to about 45\% at L5, for all
three backbones identically); blur and salt-and-pepper leave the digits
readable at every level, confirming that stroke shape, not texture, is what
MNIST classification runs on.

\textbf{Levels versus quality, in one sentence.}\quad Reading
Table~\ref{tab:phase_c2} down the level columns and
Figure~\ref{fig:c2_type_psnr} along the quality axis says the same thing two
ways: severity level is a poor common currency across degradation types,
and once converted into measured image quality, resolution is where the
accuracy actually goes.
